\subsection{Expediente verificable del perfil AES/SHA}
\label{sec:ais-lean}
El directorio \path{formal/AISProfile.lean} usa el espacio de nombres
\texttt{AISProfile}. El tipo de la clave tiene 256 bits, el del bloque
128 y el argumento completo de SHA-256 384. AES se parametriza por
cifrado y descifrado mutuamente inversos, y SHA por la función que actúa
sobre el mensaje ya serializado. La biblioteca del producto debe realizar
esas primitivas con el orden de bits y el padding de sus especificaciones
\cite{aes,sha}.

\paragraph{1. Correspondencia y siembra (.1, .2, .4).}
\proofat{AISProfile.referencePresent}{referencePresent} y \proofat{AISProfile.referenceAbsent}{referenceAbsent} comparan el relleno por la derecha del §996 y las funciones de la
tupla del §997 especializadas a la petición completa. Coinciden tanto
el bloque emitido como el estado posterior. La entrada vacía y el
bloque de 128 ceros inducen el mismo paso en este ejemplo
(\proofat{AISProfile.allZeroIsAutonomous}{allZeroIsAutonomous}). La interfaz física sólo admite un bloque
fresco. \proofat{AISProfile.requestBound}{requestBound} comprueba la longitud frente al §207.

Para nuestra instancia elegimos \(SM=\bits^{256}\), \(PS=\{\varnothing\}\) y
\(S=\bits^{256}\). Cada material \(m=B_0\|B_1\) consume dos bloques distintos
del PTG.2 interno. Las dos 4-tuplas completas, según (3.3) y (3.4), son
\[
  (SM,PS,S,\phi_{\rm seed}),\quad \phi_{\rm seed}(m,\varnothing)=m,
  \qquad
  (SM,PS,S,\phi_{\rm reseed}),\quad
  \phi_{\rm reseed}(s,m,\varnothing)=s\oplus m .
\]
El §145 contempla la copia inicial y el §148 exige usar el estado vigente
en el reseed. La traslación XOR adopta este requisito sin atribuir entropía
al estado anterior. En Lean, \proofat{AISProfile.seedReachable}{seedReachable} establece que puede
representarse cualquier clave, y \proofat{AISProfile.reseedInverse}{reseedInverse} da la inversa en
\(m\) para cada \(s\). Si el estado, la vista y el primer bloque se incluyen
en los condicionamientos apropiados, las cotas puntuales por bloque
se multiplican para dar la cota de minentropía conjunta de la
sección de arquitectura. Acreditar esas distribuciones y el origen PTG.2
sigue correspondiendo a la fuente \cite[§§143--150 y 205--207]{ais}.

\paragraph{Influencia física sobre el estado (§381).}
Para un estado \(s\), \proofat{AISProfile.influenceAtState}{influenceAtState} prueba que la transición
depende de la entrada fresca si y sólo si existen dos sufijos \(b_0,b_1\)
con \(\SHA(s\|b_0)\ne\SHA(s\|b_1)\).
Con \(s=0^{256}\), \(a_0=p\), \(a_1=p\oplus 1\), los sufijos son
los enteros 0 y 1 codificados en 16 octetos, más significativo primero.
Los estados obtenidos con SHA-256 sobre los 48 octetos son, en ese orden,
\begin{center}\small
\path|17b0761f87b081d5cf10757ccc89f12be355c70e2e29df288b65b30710dcbcd1|\\
\path|86f77354f38f799c3e3d853dec927db2d921a6d80cc206b01b9d0cf485324fce|.
\end{center}
El vector demuestra influencia en ese estado. Para reclamar dependencia
en \emph{cada} estado admitido se debe justificar que ninguna sección
\(b\mapsto\SHA(s\|b)\) es constante. Si se encuentra un estado con sección
constante, \proofat{AISProfile.constantSliceCollision}{constantSliceCollision} da dos mensajes distintos
\(s\|0, s\|1\) con el mismo digest. Esto relaciona la búsqueda de tal
estado con la resistencia a colisiones del SHA real. En el modelo de tabla
aleatoria, para cada \(s\) los tres sufijos distintos 0, 1 y 2 arrojan un
mismo digest con probabilidad \(2^{-512}\). La unión sobre \(2^{256}\)
estados acota la probabilidad de alguna sección constante por
\(2^{-256}\). El vector y la estimación de modelo tienen alcances distintos
\cite[§§381 y 996--998]{ais}.

\paragraph{2. Estado efectivo (.3).}
\proofat{AISProfile.plantInjective}{plantInjective} inyecta los estados de 246 bits en los de 256.
Para separar su comportamiento fijamos cinco textos AES distintos.
La entrada \(a=p\oplus b\) realiza cada texto \(b\) entre las elecciones
admitidas. Si dos estados tienen las mismas respuestas a todas las entradas,
sus cinco respuestas coinciden. \proofat{AISProfile.distinguishSignatures}{distinguishSignatures} y
\proofat{AISProfile.effectiveClasses}{effectiveClasses} demuestran la implicación y la inyección en las
clases de comportamiento bajo la premisa de firmas distintas. En el modelo
de permutaciones independientes del paso 3 de la sección~\ref{sec:ais-drgg-derive},
la probabilidad de que falle esta separación, incluso sobre las
\(L=2^{256}\) claves, se acota por \(\binom{L}{2}/(M)_5<2^{-128}\), donde
\(M=2^{128}\) y \((M)_5=M(M-1)\cdots(M-4)\).
\proofat{AISProfile.signatureCollisionScale}{signatureCollisionScale} comprueba la desigualdad entera.
La realización física debe sustentar el tamaño del soporte admisible. Si la
cota conjunta de la fuente es de 250,88 bits, su soporte tiene cardinalidad
superior a \(2^{246}\) (\proofat{AISProfile.seedProfileThreshold}{seedProfileThreshold} comprueba el umbral).
\proofat{AISProfile.effectiveSupported}{effectiveSupported} acepta cualquier inyección en ese soporte, que
puede diferir de la imagen de la extensión por ceros.
La separación de respuestas establece distinguibilidad extensional.
La dificultad de determinar el estado eficaz por observación, en el
sentido computacional de §§91--95, requiere además el argumento
criptográfico del §998 y el análisis de exposición de la instancia.

\paragraph{3. Derivación de predicción (.5--.8).}
Para una cuantificación computacional explícita definimos una sola
premisa local sobre las primitivas concretas. Para cada texto conocido
\(b\), los tests de coste a lo sumo \(T\) tienen ventaja de aceptación
unilateral de a lo sumo \(\alpha(T)\) entre las dos parejas
\[
  (\AES_K(b),\SHA(K\|b)),\quad (U_{128},U_{256}),\qquad K\sim U_{256},
\]
con uniformes mutuamente independientes en la segunda. El test ve
\emph{ambas} coordenadas. Esta hipótesis conjunta sobre la clave AES que
también sirve de prefijo a SHA es una condición suficiente explícita
de esta cuantificación. La PRP de AES por sí sola no establece la ley
conjunta. El §998 formula su justificación cualitativa en términos de AES
y la unidireccionalidad de la transición. El modelo de cifrado ideal y tabla
SHA independiente de la sección~\ref{sec:ais-drgg-derive} es otra forma de
explicitar los supuestos.

Sean \(R_n(h)\) el número de claves iniciales que llevan a éxito tras
\(n\) peticiones reales con historial \(h\) e \(I_n(h)\) el número de
secuencias ideales ganadoras de \(n\) bloques y una clave terminal.
El observador elige la entrada siguiente a partir de su vista y puede
comprobar el estado terminal. Si fijamos las monedas independientes del
observador, su prueba de continuación ejecuta lo que queda de la época
desde el par recibido. El coste \(T\) debe cubrir esa simulación. Con
\(d/(LM)\) como cota de la ventaja local, la sustitución de un paso da
\[
  M R_{n+1}(h)\ \le\ \sum_y R_n(hy)+d .
\]
\proofat{AISProfile.primitiveToPairBound}{primitiveToPairBound} obtiene la desigualdad aplicando la hipótesis
de un uso al test de continuación. \proofat{AISProfile.adaptiveHybrid}{adaptiveHybrid} la itera sobre
la tupla, y \proofat{AISProfile.accumulatedSlack}{accumulatedSlack} calcula la suma de los errores. Tras
normalizar los conteos y promediar sobre las monedas se obtiene
\[
  \Pr[\mathrm{acierto\ real}]\ \le\
  \Pr[\mathrm{acierto\ ideal}] + n\,d/(LM).
\]
Para un objetivo de \(t\) bits en posiciones previamente especificadas, la
vista ideal que sólo contiene coordenadas disjuntas y, en .7, la clave
terminal, deja \(2^t\) valores equiprobables. Cada conjetura por vista
acierta exactamente uno (\proofat{AISProfile.idealTargetCount}{idealTargetCount}). Las posiciones pueden
pertenecer a varios bloques. Las entradas elegidas han de depender de la
información permitida, sin recibir por otra vía los propios bits ocultos.
El caso de sucesor da .5, el de predecesor .6 y el de compromiso posterior
con la última entrada conocida .7. La elección de la entrada, incluso nula
o repetida, ya forma parte del observador en la inducción para .8.
Para \(N\) peticiones y semilla uniforme, esto da \(2^{-t}+N\alpha(T)\).
Si el estado inicial sólo dispone de una densidad puntual a lo sumo
\(2^{256-h}\) respecto de la uniforme, dada la vista previa,
\proofat{AISProfile.seedMassTransfer}{seedMassTransfer} permite multiplicar el límite de éxito por
ese factor. Se limita el resultado a uno. Para atribuir un trabajo
cercano a \(2^{\min(240,t)}\) según §214 hay que justificar \(T\),
\(N\), la semilla y el valor de \(\alpha(T)\) aplicable a las primitivas.
El kernel comprueba la reducción y su contabilidad, incluida la dependencia
en \(N\). La estimación de la dureza concreta requiere el argumento del
solicitante \cite[§§121--122, 211--214 y 998]{ais}.

\paragraph{4. Interpretación de .9 y .10.}
Fijados \(b\) y el estado divulgado \(z\), cualquier preimagen candidata
\(s'\) de la transición verifica exactamente \(\SHA(s'\|b)=z\)
(\proofat{AISProfile.transitionPreimage}{transitionPreimage}). Por ello el argumento del §998 utiliza
resistencia a obtener \emph{cualquier} preimagen válida en esta familia
de mensajes, con \(b\) conocido y la distribución de estados efectiva.
Para .7 la última entrada se conoce y permanece fijada según §213.
Para .9 explicitamos también la interfaz que acepta una preimagen con
otra entrada adicional. El teorema \proofat{AISProfile.alternativeInputPreimage}{alternativeInputPreimage} convierte
un predecesor \((s',a')\) en el mensaje \(s'\|(p\oplus\iota(a'))\),
y un mensaje válido \(s'\|b'\) en el predecesor
\((s',p\oplus b')\). Esta interfaz requiere resistencia a preimagen
para mensajes de 384 bits con ambos campos libres, sobre el desafío
que entrega la transición real. Una hipótesis sólo sobre sufijo fijado
no cubre la elección de otro sufijo \cite[§§212--215]{ais}.
Para .10 el expediente incorpora los ensayos sobre la salida real y su
representación. El §380 remite a .3 y .5--.8 para PTG.3.2.
La 9-tupla, las dos 4-tuplas y la matriz de la sección~\ref{sec:drg}
permiten ampliar el análisis a la clase DRG.3 autónoma.

\paragraph{5. Qué conserva la ITS con dos lecturas de memoria.}
\proofat{AISProfile.leftInverse}{leftInverse} y \proofat{AISProfile.rightInverse}{rightInverse} demuestran la biyección en
la entrada física para todo estado. Ahora fíjese el estado anterior \(s\),
supóngase \(A\) uniforme y divúlguese también el siguiente estado.
Con \(h_s(a)=\SHA(s\|(p\oplus a))\), \(t_z=|h_s^{-1}(z)|\) y \(M=2^{128}\),
se obtiene, sin hipótesis de dureza sobre SHA,
\[
  \dist(P_{Y,S^+\mid s}, U_{128}\otimes P_{S^+\mid s})
  =1-\sum_z (t_z/M)^2,\qquad
  \pg(Y\mid s,S^+)=|\operatorname{im}h_s|/M .
\]
La inversa recupera \(a\) a partir de \(s,y\), de modo que \(S^+\) es función de
\(s,y\) (\proofat{AISProfile.snapshotGraph}{snapshotGraph}). En cada fibra hay \(t_z\) celdas de masa
\(1/M\) y las restantes tienen masa cero. El ideal asigna \(t_z/M^2\)
a cada celda. Sumar las diferencias absolutas y dividir por dos da
\(t_z(M-t_z)/M^2\) por columna. Cada fibra ocupada aporta \(1/M\) al acierto
óptimo. Los teoremas \proofat{PTGITS.fiber_deviation}{PTGITS.fiber\_deviation},
\proofat{PTGITS.fiber_guess}{PTGITS.fiber\_guess} y \proofat{PTGITS.fiber_collision}{PTGITS.fiber\_collision} verifican
estos conteos. La fórmula expresa el resultado para las fibras efectivas
de SHA sin atribuirle inyectividad.
